\begin{theorem}[Hasse--Davenport multiplication formula]\label{mod:finitefield-hassedavenportproduct}
Let $k$ be a finite field, let $\ell\mid(\lvert k\rvert-1)$, let
$\eta\in\operatorname{FiniteMulChar}(k)$ satisfy
$\operatorname{orderOf}(\eta)=\ell$, and let
$\chi\in\operatorname{FiniteMulChar}(k)$.  If
$\psi\in\operatorname{FiniteAddChar}(k)$ is nontrivial, then
\[
 \boxed{
 \chi\!\left((\ell:k)^\ell\right)
 \tau_k(\chi^\ell,\psi)
 \prod_{j=1}^{\ell-1}\tau_k(\eta^j,\psi)
 =
 \tau_k(\chi,\psi)
 \prod_{j=1}^{\ell-1}\tau_k(\chi\eta^j,\psi).}
\]
Here $(\ell:k)$ is the natural-number cast into $k$, and every displayed
product over $1\leq j\leq\ell-1$ is Lean's
$\operatorname{Finset.Icc}(1,\ell-1)$ product.  In particular the
statement also covers $\ell=1$ and every case in which one of the
multiplicative characters is trivial.

Equivalently, for Mathlib's ordinary Gauss sum
$G_k(\alpha,\psi)$, with $G_k(1,\psi)=-1$, the internally proved formula
is
\[
 \prod_{j=0}^{\ell-1}G_k(\alpha\eta^j,\psi)
 =
 \alpha\!\left(((\ell:k)^{-1})^\ell\right)
 G_k(\alpha^\ell,\psi)
 \prod_{j=1}^{\ell-1}G_k(\eta^j,\psi).
\]
\LeanFile{LanglandsFirstMainLemma/FiniteField/HasseDavenportProduct.lean}
\lean{LanglandsFirstMainLemma.hasseDavenportProduct}
\leanok
\SourceText{Theorem thm:HD-product.}
\uses{mod:finitefield-characterapi,mod:finitefield-hdproductparametrization,mod:basic-finiteproducts,mod:finitefield-hassedavenportlift}
\FormalizationContract Mathlib multiplicative characters are used with their
canonical value zero at the field element zero.  The hypotheses in Lean are
exactly
\texttt{hdiv : ell $\mid$ Fintype.card k - 1},
\texttt{heta : orderOf eta = ell}, and
\texttt{hpsi : psi $\ne$ 1}; there is no hypothesis on $\chi$.
The multiplication theorem is not imported or assumed in any equivalent
form.  The only Davenport--Hasse input is the already proved lifting theorem
\ref{mod:finitefield-hassedavenportlift}.
\end{theorem}
\begin{proof}
First work with ordinary Gauss sums.  Every multiplicative character,
including $1$, is extended by zero at zero.  For a family
$\alpha_i\in\operatorname{FiniteMulChar}(k)$ define
\[
 J(\alpha_1,\ldots,\alpha_r)
 =\sum_{\substack{x_i\in k\\ \sum_i x_i=1}}
    \prod_i\alpha_i(x_i).
\]
Expanding the product of the Gauss sums and grouping tuples by
$a=\sum_i x_i$ proves
\[
 \prod_iG_k(\alpha_i,\psi)
 =J(\alpha_1,\ldots,\alpha_r)
   G_k\!\left(\prod_i\alpha_i,\psi\right)
\]
when the product character is nontrivial.  No quotient is taken here.
When the product character is trivial but the family is not all trivial,
the zero fiber is retained and the separate exact computation gives
\[
 \prod_iG_k(\alpha_i,\psi)
 =-\lvert k\rvert J(\alpha_1,\ldots,\alpha_r).
\]
The all-trivial family is evaluated separately by counting the tuples with
all coordinates nonzero.  Thus neither the product-trivial branch nor a
possibly vanishing Gauss sum is cancelled.

For binary Jacobi sums put
$\mathcal J(\alpha,\beta)=-J(\alpha,\beta)$.  If $K/k$ is a finite-field
extension and $\alpha\ne1$, the lift from
\ref{mod:finitefield-hassedavenportlift}, together with the preceding
Gauss--Jacobi identities, proves
\[
 \mathcal J(\alpha\circ N_{K/k},\beta\circ N_{K/k})
 =\mathcal J(\alpha,\beta)^{[K:k]}.
\]
For $\alpha\beta\ne1$, the Gauss sum cancelled in this argument is first
proved nonzero.  For $\alpha\beta=1$, use instead the exact formula
$J(\alpha,\alpha^{-1})=-\alpha(-1)$.  This is the required
product-character-trivial lifting branch, not an implicit division.

We next prove the hard Jacobi multiplication identity without assuming the
desired Gauss multiplication formula.  Choose a primitive $\ell$th root
$\zeta\in k$, whose existence follows from
$\ell\mid\lvert k^\times\rvert$, and define the multiplicative monic
polynomial weight
\[
 w(P)=\alpha\!\left(
   \prod_{j=0}^{\ell-1}P^{\mathrm{rev}}(\zeta^j)\right).
\]
For monic polynomials, reversal is multiplicative, so $w$ is a genuine
monoid homomorphism.  Fourier evaluation gives a bijection from monic
polynomials of degree $\ell-1$ to tuples
$(s_j)_{j\in\mathbf Z/\ell\mathbf Z}$ satisfying
$\sum_j s_j=(\ell:k)$.  The divisibility hypothesis also proves
$(\ell:k)\ne0$, so scaling this hyperplane to $\sum_jx_j=1$ yields
\[
 A_{\ell-1}(w)
 =\alpha((\ell:k))^\ell
   J(\underbrace{\alpha,\ldots,\alpha}_{\ell\ \mathrm{terms}}).
\]

For the degree-$r$ finite extension $k_r/k$, the closed-point power sum of
this weight is
\[
 B_r(w)=\sum_{y\in k_r}
   \alpha\!\left(N_{k_r/k}(1-y^\ell)\right).
\]
The key closed-point calculation is coefficient- and sign-exact.  If
$P=\operatorname{minpoly}_k(y)$ and $h=[k_r:k(y)]$, then for every
$z\in k$,
\[
 \bigl(P^{\mathrm{rev}}(z)\bigr)^h
 =N_{k_r/k}(1-zy).
\]
It follows from the reverse characteristic polynomial
$\det(1-Xm_y)$; hence it remains valid at every zero of either side and
introduces no sign.  Multiplying over the powers of $\zeta$ and using
$\prod_j(1-\zeta^jy)=1-y^\ell$ proves the displayed formula for $B_r$.
Minimal-polynomial fibers have their exact degree as cardinality, so the
factor $d$ and exponent $r/d$ in the closed-point definition are both
preserved.

Assume now that $\alpha^\ell\ne1$.  Finite character orthogonality for the
power map, including its zero fiber, and the signed binary Jacobi lift give
\[
 B_r(w)=(-1)^{r-1}
   \sum_{j=1}^{\ell-1}J(\alpha,\eta^j)^r
 \qquad(r>0).
\]
The integral Newton recurrence supplied by the Euler-product machinery
then gives
\[
 A_{\ell-1}(w)=
   \prod_{j=1}^{\ell-1}J(\alpha,\eta^j).
\]
Comparing the two evaluations of $A_{\ell-1}(w)$ proves the hard Jacobi
identity
\[
 \alpha((\ell:k))^\ell
 J(\alpha,\ldots,\alpha)
 =\prod_{j=1}^{\ell-1}J(\alpha,\eta^j).
\]
Combining it with the generalized and binary Gauss--Jacobi identities
proves the ordinary multiplication formula.  The generalized Jacobi sum
is cancelled only after its nonvanishing follows from
$G_k(\alpha,\psi)^\ell
=J(\alpha,\ldots,\alpha)G_k(\alpha^\ell,\psi)$ and the nonvanishing of
both Gauss sums.

It remains to treat $\alpha^\ell=1$.  The finite character group is cyclic,
and exact order of $\eta$ parametrizes every solution of
$\gamma^\ell=1$ as $\eta^i$.  Thus $\alpha=\eta^i$, and cyclic reindexing
of $j\mapsto i+j$ identifies the two complete Gauss products.  Moreover
$\alpha(((\ell:k)^{-1})^\ell)=1$ and
$G_k(\alpha^\ell,\psi)=G_k(1,\psi)=-1$, giving the ordinary formula with
the correct sign.  This branch contains $\alpha=1$ and all exceptional
trivial-character factors.

Finally apply the ordinary identity to $\alpha=\chi^{-1}$ and
$\eta^{-1}$.  Inversion permutes the nonzero powers of $\eta$, and
\[
 \tau_k(\theta,\psi)=-G_k(\theta^{-1},\psi),
 \qquad
 \chi^{-1}\!\left(((\ell:k)^{-1})^\ell\right)
 =\chi\!\left((\ell:k)^\ell\right).
\]
Each side contains exactly $\ell$ ordinary Gauss sums, so all minus signs
cancel.  Reindexing the finite products gives precisely the boxed
Langlands-normalized identity.
\end{proof}
